% ============================================================
%  Chapter 2, Related Work
% ============================================================
\chapter{Related Work}
\label{ch:irodalom}

\section{The Kalman Filter and Its Extensions}
\label{sec:kalman}

The Kalman filter~\cite{kalman1960new} is the minimum-variance, unbiased, linear
estimator for systems whose state $x_t \in \mathbb{R}^n$ evolves according to
a linear dynamical model perturbed by Gaussian noise:
\begin{align}
    x_t &= F x_{t-1} + w_t, & w_t &\sim \mathcal{N}(0, Q), \label{eq:kf_process}\\
    y_t &= H x_t   + v_t,   & v_t &\sim \mathcal{N}(0, R), \label{eq:kf_measure}
\end{align}
where $F$ is the state-transition matrix, $H$ the observation matrix,
$Q \in \mathbb{R}^{n \times n}$ the process-noise covariance, and
$R \in \mathbb{R}^{m \times m}$ the measurement-noise covariance. The filter
operates in two alternating half-steps.
In the prediction step (time update), before the measurement $y_t$ arrives,
the filter propagates the state estimate and the error covariance forward
in time:
\begin{align}
    \hat{x}_{t|t-1} &= F \hat{x}_{t-1|t-1}, \label{eq:kf_predict_x}\\
    P_{t|t-1}       &= F P_{t-1|t-1} F^\top + Q, \label{eq:kf_predict_P}
\end{align}
where $P_{t|t-1} \in \mathbb{R}^{n \times n}$ is the predicted error
covariance. A large $P_{t|t-1}$ indicates high prediction uncertainty;
a large $Q$ means the system state changes rapidly and unpredictably.

In the correction step (measurement update), once the observation $y_t$
arrives, the filter corrects its estimate by the
\emph{innovation} $e_t = y_t - H\hat{x}_{t|t-1}$, weighted by the Kalman
gain $K_t$:
\begin{align}
    K_t           &= P_{t|t-1} H^\top \bigl(H P_{t|t-1} H^\top + R\bigr)^{-1},
                     \label{eq:kalman_gain}\\
    \hat{x}_{t|t} &= \hat{x}_{t|t-1} + K_t e_t, \label{eq:kalman_update}\\
    P_{t|t}       &= (I - K_t H)\,P_{t|t-1}. \label{eq:kf_update_P}
\end{align}
Equation~\eqref{eq:kalman_gain} encodes the fundamental noise-balance
intuition. When the measurement noise $R$ is large relative to the predicted
covariance $H P_{t|t-1} H^\top$, the denominator dominates and $K_t \to 0$:
the filter distrusts the observation and relies on its prior state. When $R
\to 0$, $K_t \to H^{-1}$ and the observation nearly overwrites the state
entirely.

For the scalar case ($F = H = 1$), the gain simplifies to
\begin{equation}
    K_t = \frac{P_{t|t-1}}{P_{t|t-1} + R}.
    \label{eq:kalman_gain_scalar}
\end{equation}
Under steady-state conditions the error covariance $P_{t|t-1}$ converges to a
fixed value $\bar{P}$, and the steady-state gain depends only on the ratio
$Q/R$: a high process noise relative to measurement noise drives $K$
towards~1 (trust the measurement), while a low $Q/R$ ratio drives $K$
towards~0 (trust the model). This $Q/(Q+R)$ ratio is the structural basis
for the KCMamba gain computation (Section~\ref{sec:kalman_gain}).

\subsection{Learned Kalman Filters}

In the neural Kalman literature, the matrices $F, H, Q, R$ are parameterised
by neural networks~\cite{krishnan2015deep,fraccaro2017disentangled}.
The \emph{Deep Kalman Filter}~\cite{krishnan2015deep} allows non-linear
transitions, and the \emph{Kalman Variational Autoencoder}~\cite{fraccaro2017disentangled}
learns disentangled representations of latent states.
KCMamba builds on these ideas but operates as a sequence forecaster rather
than as an autoencoder.

\section{State Space Models}
\label{sec:ssm}

\subsection{S4}

S4~\cite{gu2021s4} reformulates a sequence model as a parameterised ordinary
differential equation (ODE) over a latent state $h(t) \in \mathbb{R}^N$:
\begin{align}
    h'(t) &= A h(t) + B u(t), \label{eq:s4_continuous}\\
    y(t)  &= C h(t). \label{eq:s4_output}
\end{align}
Here $u(t)$ is the scalar input, $A \in \mathbb{R}^{N \times N}$ is the
state matrix, and $B, C \in \mathbb{R}^N$ are the input and output projections.
The critical insight of Gu et al.\ is that $A$ can be initialised as the
\emph{HiPPO} (High-order Polynomial Projection Operator) matrix, which is
designed to compress the history of $u$ into the current state $h(t)$ with
optimal polynomial approximation. This initialisation avoids the
vanishing-gradient problem that plagues vanilla RNNs over long sequences,
since the HiPPO structure preserves information from the distant past without
exponential decay.

To apply the model to discrete sequences the continuous system is discretised
with step size $\Delta$ using the zero-order hold (ZOH) rule:
\begin{align}
    h_t &= \bar{A} h_{t-1} + \bar{B} u_t, \label{eq:s4_discrete_h}\\
    y_t &= C h_t, \label{eq:s4_discrete_y}
\end{align}
where $\bar{A} = e^{\Delta A}$ and $\bar{B} = (e^{\Delta A} - I)A^{-1}B$.
Because $\bar{A}$ is diagonal in the frequency domain the entire convolution
$y = \bar{C} * u$ can be computed as a single element-wise product after an
FFT, reducing sequence-level complexity to $O(T \log T)$.

A fundamental limitation of S4 is that $A$, $B$, $C$, and $\Delta$ are
\emph{input-independent}: the same learned parameters are applied at every
time step regardless of the input content. The model therefore cannot adapt
its memory or forgetting behaviour in response to what it actually observes
in the current sequence.

\subsection{Mamba}

Mamba~\cite{gu2023mamba} addresses this limitation directly by making the
state-space parameters \emph{input-dependent}. The matrices $B$, $C$, and the
time step $\Delta$ are no longer fixed scalars but functions of the current
input $u_t$:
\begin{align}
    B_t      &= \mathrm{Linear}(u_t),  \\
    C_t      &= \mathrm{Linear}(u_t),  \\
    \Delta_t &= \mathrm{softplus}\bigl(\mathrm{Linear}(u_t) + \delta_\mathrm{bias}\bigr).
\end{align}
The time step $\Delta_t$ determines the discrete transition matrices
$\bar{A}_t = e^{\Delta_t A}$, $\bar{B}_t = (\bar{A}_t - I) A^{-1} B_t$.
This mechanism, called the \emph{Selective State Space} (S6) layer, lets the
model selectively remember or forget content at each step. When
$\Delta_t \to 0$ the state barely changes ($\bar{A}_t \approx I$),
preserving memory. When $\Delta_t$ is large the state is updated strongly,
discarding older information.

The input-dependent parameters break the global convolution trick that made
S4 efficient, because there is no longer a single kernel to FFT. Gu and Dao
solve this with a \emph{hardware-efficient parallel scan} (the Blelloch
prefix-scan algorithm), applied on the GPU's SRAM to avoid expensive
high-bandwidth memory round-trips. The scan runs in $O(\log T)$
synchronisation steps on $T$ parallel threads, achieving $O(T)$ total
work and $O(T)$ memory while avoiding the sequential bottleneck of recurrent
models.

\subsection{Fair Mamba}

The comparison baseline, referred to as ``Fair Mamba'', is a channel-wise Mamba
with identical hyperparameters and training protocol as KCMamba. The ``fair''
label means the parameter budget and data are identical, so any difference in
results comes solely from the architecture.

\section{Long-Term Time Series Forecasting}
\label{sec:ltsf}

\subsection{Problem Formulation}

Long-term time series forecasting (LTSF) considers an input window
$\mathbf{X} = (x_1, \ldots, x_L) \in \mathbb{R}^{L \times d}$ of $L$
historical steps across $d$ channels, and requires producing a prediction
$\hat{\mathbf{X}} = (\hat{x}_{L+1}, \ldots, \hat{x}_{L+H}) \in \mathbb{R}^{H \times d}$
for the next $H$ steps. The objective is to minimise the normalised mean
squared error
\begin{equation}
    \mathrm{MSE} = \frac{1}{H d}
        \sum_{h=1}^{H} \sum_{j=1}^{d}
        \frac{(x_{L+h,j} - \hat{x}_{L+h,j})^2}{\mathrm{Var}(x_{\cdot,j})},
    \label{eq:mse}
\end{equation}
where normalisation by the per-channel variance makes results comparable
across datasets with different scales. The \emph{long-term} qualifier refers
to the common regime $H \geq L$ or $H$ comparable to $L$: the model must
extrapolate well beyond the lookback window, which is substantially harder
than short-term forecasting where $H \ll L$.

\subsection{Reference Models}

\emph{Informer}~\cite{zhou2021informer} was the first Transformer variant
specifically designed for long-term forecasting. Standard self-attention has
$O(L^2)$ time and memory complexity, which becomes prohibitive for long input
windows. Informer addresses this with \emph{ProbSparse attention}: it
identifies the $O(\log L)$ most informative query positions via a KL-divergence
estimate, reducing complexity to $O(L \log L)$. A distilling operation
halves the sequence length between encoder layers so the model can handle
lookback windows of several thousand steps.

\emph{Autoformer}~\cite{wu2021autoformer} decomposes each series into a
trend component (estimated with a moving-average kernel) and a seasonal
residual, applying the Transformer only to the residuals. Rather than
dot-product attention it uses an \emph{Auto-Correlation} mechanism: it
identifies the lag $\tau$ maximising the autocorrelation of the query
signal and aggregates values at those lag offsets. The resulting mechanism
is $O(L \log L)$ via FFT and provides an inductive bias suited to periodic
data.

\emph{DLinear}~\cite{zeng2023transformers} discards all of this complexity
and applies a single linear layer to each lookback window to produce the
forecast, optionally after trend-seasonal decomposition. Despite its
simplicity, DLinear matched or outperformed many Transformer variants on
standard benchmarks, demonstrating that architectural complexity is not
always necessary for good forecasting performance. It remains a relevant
baseline because any new model must justify its additional overhead relative
to this linear reference.

\subsection{Datasets}
\label{sec:datasets}

The Exchange-Rate dataset~\cite{lai2018modeling} contains eight foreign exchange
rates sampled daily from 1990 to 2016, giving 7\,588 time steps per channel.
Exchange-rate data has near-unit-root dynamics: each daily increment is
close to a random walk with very little mean reversion. This makes it a
challenging benchmark for autoregressive models, and a natural setting
for evaluating the Kalman structure where the signal-to-noise ratio is
intrinsically low.

ETTh1 and ETTh2~\cite{zhou2021informer} are both collected from
Chinese transformer stations, recording hourly electricity load and oil
temperature for two years (17\,420 steps, $d=7$ channels). ETTh1 has a
clear daily periodicity, while ETTh2 is noisier and more irregular. The two
datasets provide contrasting difficulty levels for noise-robustness
evaluation.

\section{Noise Robustness}
\label{sec:robustness}

The noise sensitivity of time series models is rarely studied in
a systematic way. PatchTST~\cite{nie2022patchtst} indirectly improves noise
tolerance through patch tokenisation, but robustness is not its primary goal.

In this work the evaluation protocol is:
$\tilde{x}_t = x_t + \varepsilon_t$,
$\varepsilon_t \sim \mathcal{N}(0, \sigma^2 I)$, $\sigma \in \{0, 1, 3, 5\}$.
Both training and test sets are corrupted with the same $\sigma$, so this
is uniform data-quality degradation, not a domain-shift scenario.
